% ═══════════════════════════════════════════════════════════════
% ML-01b: Musterlösung - Sich vorstellen + Verb ser
% Abgeleitet aus LE-01b (Abschnitt 9: Lösungen)
% ═══════════════════════════════════════════════════════════════

\documentclass[11pt, a4paper]{article}

\usepackage[utf8]{inputenc}
\usepackage[T1]{fontenc}
\usepackage[ngerman]{babel}
\usepackage{helvet}
\renewcommand{\familydefault}{\sfdefault}

\usepackage[margin=2cm]{geometry}
\usepackage{enumitem}
\usepackage{tabularx}
\usepackage{booktabs}

\usepackage{xcolor}
\usepackage{tcolorbox}

\usepackage{fancyhdr}
\usepackage{lastpage}

% Farben
\definecolor{spanisch}{HTML}{c41e3a}
\definecolor{grau}{HTML}{888888}
\definecolor{hellgrau}{HTML}{f5f5f5}
\definecolor{orange}{HTML}{b35c00}
\definecolor{loesung}{HTML}{c62828}

\newcommand{\fachfarbe}{spanisch}

% Variablen
\newcommand{\thema}{Sich vorstellen + Verb ser}
\newcommand{\sequenz}{01}
\newcommand{\block}{b}
\newcommand{\fach}{Spanisch}
\newcommand{\klasse}{7}
\newcommand{\maxpunkte}{20}

% Kopf-/Fußzeile
\pagestyle{fancy}
\fancyhf{}
\renewcommand{\headrulewidth}{0.4pt}
\renewcommand{\footrulewidth}{0.4pt}
\lhead{\fach{} -- Klasse \klasse}
\chead{\textcolor{loesung}{\textbf{ML-\sequenz\block{} (MUSTERLÖSUNG)}}}
\rhead{}
\lfoot{}
\rfoot{Seite \thepage\ von \pageref{LastPage}}

% Befehle
\newcommand{\punktebox}[1]{\hfill\fbox{\small \_/#1}}
\newcommand{\loesung}[1]{\textcolor{loesung}{\textbf{#1}}}
\newcommand{\luecke}[1]{\rule{#1}{0.4pt}}
\newcommand{\es}[1]{\textcolor{\fachfarbe}{\textbf{#1}}}

% Merkkasten
\newtcolorbox{merkkasten}[1][]{
    colback=hellgrau,
    colframe=\fachfarbe,
    boxrule=1pt,
    arc=0pt,
    left=8pt, right=8pt, top=8pt, bottom=8pt,
    fonttitle=\bfseries\color{\fachfarbe},
    title=#1
}

% Hinweis
\newtcolorbox{hinweis}{
    colback=white,
    colframe=orange,
    boxrule=1pt,
    arc=0pt,
    left=5pt, right=5pt, top=5pt, bottom=5pt,
    fonttitle=\bfseries,
    title=Hinweis
}

% Aufgabe
\newenvironment{aufgabe}[1]{%
    \vspace{10pt}
    \noindent\textbf{\textcolor{\fachfarbe}{Aufgabe #1}}
    \vspace{5pt}
    \begin{list}{}{\leftmargin=0pt}
    \item[]
}{%
    \end{list}
}

% ═══════════════════════════════════════════════════════════════
\begin{document}

% Kopfbereich
\begin{center}
    {\Large\bfseries\textcolor{\fachfarbe}{Arbeitsblatt: \thema}}\\[0.3em]
    {\large\textcolor{loesung}{\textbf{--- MUSTERLÖSUNG ---}}}
\end{center}

\vspace{5pt}

\noindent
\begin{tabularx}{\textwidth}{|X|X|X|}
\hline
\textbf{Name:} \textcolor{loesung}{Musterschüler/in} &
\textbf{Klasse:} \klasse &
\textbf{Datum:} \textcolor{loesung}{XX.XX.XXXX} \\
\hline
\end{tabularx}

\vspace{15pt}

% ═══════════════════════════════════════════════════════════════
% MERKKASTEN (mit Lösungen)
% ═══════════════════════════════════════════════════════════════

\begin{merkkasten}[Merkkasten: Das Verb \textit{ser} (sein)]

\begin{center}
\begin{tabular}{|l|l|l|}
\hline
\textbf{Person} & \textbf{Pronomen} & \textbf{ser} \\
\hline
1. Sg. & yo (ich) & \loesung{soy} \\
\hline
2. Sg. & tú (du) & \loesung{eres} \\
\hline
3. Sg. & él/ella (er/sie) & \loesung{es} \\
\hline
1. Pl. & nosotros/-as (wir) & \loesung{somos} \\
\hline
\end{tabular}
\end{center}

\end{merkkasten}

\vspace{10pt}

% ═══════════════════════════════════════════════════════════════
% AUFGABE 1 (mit Lösungen)
% ═══════════════════════════════════════════════════════════════

\begin{aufgabe}{1}
Verbinde die Pronomen mit der richtigen Form von \textit{ser}. \punktebox{4}
\end{aufgabe}

\begin{center}
\begin{tabular}{rcl}
yo & \loesung{------} & \loesung{soy} \\[0.8em]
tú & \loesung{------} & \loesung{eres} \\[0.8em]
él/ella & \loesung{------} & \loesung{es} \\[0.8em]
nosotros & \loesung{------} & \loesung{somos} \\
\end{tabular}
\end{center}

\textit{\textcolor{grau}{Je 1P pro richtige Zuordnung}}

\vspace{15pt}

% ═══════════════════════════════════════════════════════════════
% AUFGABE 2 (mit Lösungen)
% ═══════════════════════════════════════════════════════════════

\begin{aufgabe}{2}
Setze die richtige Form von \textit{ser} ein. \punktebox{4}
\end{aufgabe}

\begin{enumerate}[label=\alph*)]
    \item Me llamo María. \loesung{Soy} de Alemania.
    \item ¿Cómo te llamas? -- \loesung{Soy} Pedro.
    \item Él \loesung{es} de España.
    \item Nosotros \loesung{somos} de Rostock.
\end{enumerate}

\textit{\textcolor{grau}{Je 1P pro richtige Form}}

\vspace{15pt}

% ═══════════════════════════════════════════════════════════════
% AUFGABE 3 (mit Lösungen)
% ═══════════════════════════════════════════════════════════════

\begin{aufgabe}{3}
Vervollständige den Dialog. Nutze die Redemittel aus dem Kasten. \punktebox{6}
\end{aufgabe}

\noindent
\textbf{A:} \es{¡Hola!} \loesung{¿Cómo te llamas?}

\textbf{B:} \loesung{Me llamo [Name].} \es{¿Y tú?}

\textbf{A:} \es{Soy Carlos.} \loesung{¿De dónde eres?}

\textbf{B:} \loesung{Soy de Alemania. / Soy de [Ort].}

\textbf{A:} \es{¡Mucho gusto!}

\textbf{B:} \loesung{¡Igualmente! / ¡Mucho gusto!}

\vspace{0.5em}
\textit{\textcolor{grau}{Je 1P pro sinnvolle Ergänzung, max. 6P}}

\vspace{15pt}

% ═══════════════════════════════════════════════════════════════
% AUFGABE 4 (Beispiellösung)
% ═══════════════════════════════════════════════════════════════

\begin{aufgabe}{4}
Schreibe einen eigenen Vorstellungsdialog (mind. 6 Zeilen). \punktebox{6}
\end{aufgabe}

\textbf{Beispiellösung:}

\loesung{A: ¡Hola! ¿Qué tal?}

\loesung{B: Bien, gracias. ¿Y tú?}

\loesung{A: Muy bien. ¿Cómo te llamas?}

\loesung{B: Me llamo Ana. ¿De dónde eres?}

\loesung{A: Soy de Rostock, de Alemania.}

\loesung{B: ¡Mucho gusto! ¡Hasta luego!}

\vspace{0.5em}
\textit{\textcolor{grau}{Bewertung: Begrüßung (1P), Name (1P), Herkunft (2P), Verabschiedung (1P), Korrektheit (1P)}}

\vspace{0.5em}
\textbf{Bonus (tú/usted):} \loesung{Tú verwendet man unter Freunden und Gleichaltrigen. Usted ist die höfliche Form für Respektspersonen oder im Geschäftsleben.}

\vspace{20pt}
\hrule
\vspace{10pt}

\begin{flushright}
\textbf{Punkte:} \loesung{\maxpunkte} / \maxpunkte
\end{flushright}

\end{document}
